\subsubsection{\textbf{算法概述、流程与关键参数}}

本节在 2.1 给出的计算域 $\Omega$、状态函数 $S(\textbf x)$、孔隙率目标约束 $|\varepsilon - \varepsilon_{\textbf{target}}| \le \delta_{\varepsilon}$ 的基础上，构建一种面向 GDL 的分层随机纤维生成算法，用于生成纤维相体素集合 $V_{\textbf{fiber}}$。算法以“逐根生成——判定——写入”的迭代方式进行：每次迭代采样候选纤维的中心位置，方向与长度并进行边界裁剪；允许重叠时直接体素化写入，禁止重叠时先进行碰撞检测，通过后写入；每次写入后更新当前孔隙率 $\varepsilon$，当满足目标约束或达到最大尝试次数时终止。
具体的算法流程如下所示：

% 将流程图以内嵌方式显示（非浮动），避免被 LaTeX 浮动机制推到别处或造成两栏不平衡的大空白 
\begin{center} % 将流程图限制在单栏宽度内并在必要时缩放 
  \begin{adjustbox}{max width=\columnwidth,center} 
    \begin{tikzpicture}[ 
      node distance=5mm and 10mm,
      font=\scriptsize, 
      startstop/.style={rectangle, rounded corners, draw, fill=white, align=center, text width=26mm, inner sep=3pt, minimum height=6mm}, 
      process/.style={rectangle, draw, fill=white, align=center, text width=30mm, inner sep=3pt, minimum height=6mm}, decision/.style={diamond, aspect=2.1, draw, fill=white, align=center, text width=28mm, inner sep=1.5pt}, arrow/.style={-Latex, line width=0.6pt} ] 
      \node (start) [startstop] {开始}; 
      \node (init) [process, below=of start] {初始化：$S(x)=0$，纤维列表为空\\计算当前孔隙率 $\varepsilon$}; 
      \node (stopq) [decision, below=of init, yshift=-2mm] {是否满足停止条件？\\$|\varepsilon-\varepsilon_{\text{target}}|\le\delta_\varepsilon$\\或 tries 达上限}; 
      \node (out) [startstop, right=of stopq, xshift=12mm] {输出：$S(x)$ 与纤维列表}; 
      \node (layer) [process, below=of stopq, yshift=-2mm] {选择层 $m$ 并采样中心 $c$\\（分层 + jitter）}; 
      \node (ori) [process, below=of layer] {采样取向 $(\alpha,\beta)$\\得到方向向量 $u$}; 
      \node (len) [process, below=of ori] {采样长度 $l\in[l_{\min},l_{\max}]$\\构造端点 $a,b$}; 
      \node (clip) [process, below=of len] {边界裁剪（clipped）\\得到域内线段 $[a,b]$}; 
      \node (over) [decision, below=of clip, yshift=-2mm] {是否允许重叠？}; 
      \node (vox) [process, left=of over, xshift=-12mm] {体素化写入 $S(x)$\\记录纤维信息}; 
      \node (coll) [process, right=of over, xshift=12mm] {碰撞检测（不重叠）\\分箱/AABB 预筛\\线段最小距离 $\ge 2r$}; 
      \node (pass) [decision, below=of coll, yshift=-2mm] {检测通过？}; 
      \node (adapt) [process, below=of pass] {失败计数+1\\局部 nudge / 自适应缩短 $l_{\max}$\\缓解 jamming}; 
      \node (upd) [process, above=of vox] {更新孔隙率 $\varepsilon$}; 
      \draw [arrow] (start) -- (init); 
      \draw [arrow] (init) -- (stopq); 
      \draw [arrow] (stopq) -- node[above] {是} (out); 
      \draw [arrow] (stopq) -- node[left] {否} (layer); 
      \draw [arrow] (layer) -- (ori); 
      \draw [arrow] (ori) -- (len); 
      \draw [arrow] (len) -- (clip); 
      \draw [arrow] (clip) -- (over); 
      \draw [arrow] (over) -- node[above] {允许} (vox); 
      \draw [arrow] (over) -- node[above] {不允许} (coll); 
      \draw [arrow] (coll) -- (pass); 
      % 从 pass.west 出发：先向左，再向上，汇入 over->vox 的“允许”箭头（保持黑色、无箭头）
      \coordinate (pass_left) at ($(pass.west)+(-72mm,0)$);
      \coordinate (allow_mid) at ($(over.west)!0.5!(vox.east)$);
      \coordinate (pass_left_up) at ($(pass_left |- allow_mid)$);
      \draw [line width=0.6pt] (pass.west) -- node[below,pos=0.15]{是} (pass_left) -- (pass_left_up) -- (allow_mid);
      \draw [arrow] (pass) -- node[left] {否} (adapt); 
      % 从 adapt.east 出发：先向右一段，再向上对齐 stopq.south 的纵坐标，最后水平连到 stopq.south（改为右侧回路）
      \coordinate (adapt_right) at ($(adapt.east)+(12mm,0)$);
      \coordinate (adapt_right_up) at ($(adapt_right |- stopq.south)$);
      \draw [arrow] (adapt.east) -- (adapt_right) -- (adapt_right_up) -- (stopq.south);
      \draw [arrow] (vox) -- (upd); 
      \draw [arrow] (upd.north) |- (stopq.west); 
    \end{tikzpicture} 
  \end{adjustbox} 
  \captionof{figure}{分层随机纤维生成流程（含允许/不允许重叠分支与自适应防 jamming 策略）。} 
  \label{fig:fiber_flowchart} 
\end{center}
分层随机纤维生成算法的主要参数如下：

\begin{table}[htbp]
\centering
\begin{adjustbox}{max width=\columnwidth,center}
\begin{tabular}{lll}
\hline
\textbf{符号/名称} & \textbf{含义} & \textbf{典型取值/范围} \\
\hline
$I,J,K$ & 体素域尺寸（第 2.1.1 节） & $155\times 300\times 300$ \\
$r$ & 纤维半径（体素） & 例如 $r=3.5$ \\
$l_{\min},l_{\max}$ & 纤维长度范围（体素） & 例如 $[140,320]$ \\
$n_{\mathrm{layer}}$ & 厚度方向分层数 & 例如 $6\sim 11$ \\
$\sigma_z$（jitter） & 层内厚度方向中心抖动 & 例如 $2\sim 5$ 体素 \\
$\mu_{\alpha}^{(m)}$ & 第 $m$ 层的平面内平均主方向 & 给定/缓慢变化 \\
$\sigma_\alpha$ & 平面内角度离散度（层内“凌乱”程度） & 例如 $40^\circ\sim 80^\circ$ \\
$\sigma_\beta,\beta_{\max}$ & 离层倾角控制（贴层程度） & 例如 $\sigma_\beta=2^\circ\sim 5^\circ$ \\
\texttt{overlap} & 是否允许重叠（Y/N） & Y / N \\
\texttt{fiber\_end\_mode} & 端点边界处理方式 & clipped \\
$\varepsilon_{\mathrm{target}}$ & 目标孔隙率（第 2.1.3 节） & 例如 $0.70\sim 0.85$ \\
$\delta_\varepsilon$ & 孔隙率容差（停止阈值） & 例如 $10^{-4}\sim 10^{-3}$ \\
\texttt{max\_tries} & 最大尝试次数（防止无限循环） & 依实现设定 \\
\hline
\end{tabular}
\end{adjustbox}
\caption{分层随机纤维生成算法的主要参数。}\label{tab:fiber_params_cn}
\end{table}

\subsubsection{\textbf{厚度方向分层建模与层内位置采样}}

GDL 在厚度方向通常呈现\textbf{厚度方向层状堆叠}、\textbf{层内位置随机、乱中有序}的特征，本节主要通过分层建模与层内位置采样，来给每根纤维分配一个合理的厚度位置，并给它采样一个中心点 $\textbf c = (x,y,z)$。

取计算域中最短的维度（通常是 $I$ 轴，对应实现中的 TP\_AXIS）作为厚度方向（through-plane），其余两维为面内方向（in-plane）（对应实现中的 IP1\_AXIS 和 IP2\_AXIS）。设厚度方向的长度为 $K_{\textbf{TP}}$，为贴近实际 GDL 的结构，我们将厚度方向划分为 $n_{\mathrm{layer}}$ 个连续的层（layer），那么每一层的厚度（层厚）为
\[
h = \dfrac{K_{\textbf{TP}}}{n_{\mathrm{layer}}}
\]
其中第 $m$ 层的中心位置为
\[
z_m = \left(m + \dfrac 1 2\right)h, \quad m=0,1,\ldots,n_{\mathrm{layer}}-1
\]

考虑如何对于每根碳棒纤维选择一个层。对于\textbf{允许重叠}的情况，我们直接等概率的选取一个层 $m$；而对于\textbf{禁止重叠}的情况，为了缓解后续可能出现的 jamming（即后续纤维难以找到合适位置写入），我们采用一种基于当前层占用情况的\textbf{动态加权随机选择}策略：为了避免过早地将某些层填满导致后续纤维难以写入，我们应当让当前占用较少的层有更高的概率被选中，具体地，我们定义
\[
w_i = \dfrac{1}{cnt_i + 1}, \quad i=0,1,\ldots,n_{\mathrm{layer}}-1
\]
其中 $cnt_i$ 表示第 $i$ 层当前已写入的纤维数量，$w_i$ 则是第 $i$ 层被选中的权重。
通过这种方式，不仅可以在禁止重叠的情况下更均匀地分布纤维，避免过早地填满某些层导致后续纤维难以写入，还可以保留一定的随机性，使得生成的结构更接近实际 GDL 的“乱中有序”特征。

在前文中，我们已经找到了每一层的中心位置 $z_m$，接下来我们需要为每根纤维采样一个具体的厚度位置 $z$。为了模拟实际 GDL 中纤维在层中心较为集中、向上下逐渐稀疏的分布特征，我们在层中心位置 $z_m$ 的基础上，添加一个服从零均值正态分布的高斯抖动（jitter）：
\[
z \sim \mathcal{N}(z_m,\sigma^2)
\]
其中 $z$ 是最终的厚度位置，$\sigma$ 作为标准差则控制了层内厚度方向的离散程度（即“凌乱”程度）。根据正态分布的特点，最终大多数点会落在 $z_m$ 附近，距离中心越远概率快速变小；较小的 $\sigma$ 会使得纤维更集中在层中心，而较大的 $\sigma$ 则会使得纤维在厚度方向上分布得更为分散。通过调整 $\sigma$ 的值，我们可以模拟不同程度的层内“凌乱”，以更好地匹配实际 GDL 的结构特征。

需要注意的是，正态分布的支持是整个实数轴，因此理论上可能会采样到超出计算域厚度范围的 $z$ 值。为了确保生成的纤维能够正确地写入计算域，我们需要对采样得到的 $z$ 进行边界裁剪（clipping），将其限制在合理的范围内。

接下来，我们需要为每根纤维采样一个面内位置 $(x,y)$，为了满足 GDL 层内“凌乱”的特性，我们可以直接在计算域的 $xy$ 平面上进行均匀随机采样：
\[  
\begin{cases}
x \sim \mathcal{U}(0, I_{\textbf{IP1}}) \\
y \sim \mathcal{U}(0, I_{\textbf{IP2}})
\end{cases}
\]
其中 $I_{\textbf{IP1}}$ 和 $I_{\textbf{IP2}}$ 分别表示计算域在面内两个方向上的尺寸。通过这种方式，能够确保每一层内的纤维在面内位置上是随机分布的。据此，我们就得到了每根纤维的中心点坐标 $\textbf c = (x,y,z)$。

通过上述算法，我们成功地为每根纤维分配了一个合理的中心点坐标，并且模拟了实际 GDL 中纤维在厚度方向层状堆叠、层内位置随机但又有一定规律的分布特征。这为后续的纤维生成与写入奠定了坚实的基础。

\subsubsection{\textbf{纤维方向采样与空间位置确定}}

在上一节的内容中，我们已经为每根纤维分配了一个合理的中心点坐标。那么为了唯一确定一根三维空间中的碳棒纤维，我们还需要指定它的方向。

我们用两个角度来参数化纤维的方向：
\begin{itemize}
\item $\alpha$：纤维在面内（in-plane azimuth）的主方向角，定义为纤维轴线在 $xy$ 平面上的投影与 $x$ 轴的夹角，取值范围为 $\alpha \in [0,2 \pi)$。
\item $\beta$：纤维的离层倾角（out-of-plane tilt），定义为纤维轴线与厚度方向（$z$ 轴）的夹角。
\end{itemize}
则纤维的方向向量可定义为：
\[
\textbf u = (\cos\alpha \cos\beta, \cos\beta \sin\alpha, \sin\beta)
\]

记第 $m$ 层的平均主方向为 $\mu_{\alpha}^{(m)}$，我们在层内采样 $\alpha$ 时，以 $\mu_{\alpha}^{(m)}$ 为均值，使用一个适当的标准差 $\sigma_\alpha$ 来控制层内纤维方向的离散程度（即“凌乱”程度），从而使得每层存在一个明显的以 $\mu_{\alpha}^{(m)}$ 为主方向，纤维倾向于围绕它分布的特征。具体地，我们可以使用一个截断高斯 + 环上高斯来采样 $\alpha$：
首先采样一个“未取模”的角度：
\[
\tilde\alpha \sim \mathcal{N}(\mu_{\alpha}^{(m)}, \sigma_\alpha^2)
\]
再把它映射到 $[0,2\pi)$ 上：
\[
\alpha = \tilde\alpha \bmod 2\pi
\]
同时，为了防止高斯长尾导致过于离谱的角度，规定最大偏转角 $\alpha_{\max}$，在采样 $\tilde\alpha$ 时进行截断，限制其在 $[\mu_{\alpha}^{(m)} - \alpha_{\max}, \mu_{\alpha}^{(m)} + \alpha_{\max}]$ 范围内。
对于 $\mu_{\alpha}^{(m)}$ 的设置，可以让所有层共享一个全局主方向 $\mu_{\alpha}$；也可以让每层主方向 $\mu_{\alpha}^{(m)}$ 随层数 $m$ 缓慢变化，即 $\mu_{\alpha}^{(m)} = \mu_0 + \Delta\mu \cdot m$，以模拟实际 GDL 中不同层主方向逐渐变化的特征。

对于离层倾角 $\beta$ 的采样，我们同样使用一个截断高斯来控制纤维的贴层程度（即与厚度方向的夹角）。由于 GDL 的真实结构中纤维大多“躺在层平面内”，所以设定 $\beta$ 的均值为 $0^\circ$，表示纤维倾向于贴近层面。即 $\beta$ 满足正态分布
\[
\tlide\beta \sim \mathcal{N}(0, \sigma_\beta^2), \quad \beta = \text{clip}(\tilde\beta, -\beta_{\max}, \beta_{\max})
\]
其中 $\beta_{\max}$ 则限制了纤维的最大离层倾角，相当于是一个硬性约束。假设一根纤维的中心轴线长度为 $L$，与层平面的夹角为 $\beta$，那么应当满足
\[
L \sin|\beta| \le K_{\textbf{TP}}
\]
因为对任意 $L \ge L_{\min}$，所以要使得上式满足对于所有可能的纤维长度，则需要满足
\[
L_{\min} \sin\beta \le K_{\textbf{TP}}
\]
化简得
\[
|\beta| \le \sin^{-1}\left(\dfrac{K_{\textbf{TP}}}{L_{\min}}\right)
\]
于是定义
\[
\beta_{\max} = \sin^{-1}\left(\dfrac{K_{\textbf{TP}}}{L_{\min}}\right)
\]

由上述分布得到的方向向量 $\textbf u$，结合前面采样得到的中心点 $\textbf c$，就可以唯一确定一根纤维的空间位置和方向。

\subsubsection{\textbf{纤维长度与端点处理}}

在前面的内容中，我们已经通过中心点 $\textbf c$ 和方向向量 $\textbf u$ 确定了纤维的空间位置和方向。接下来，我们需要为每根纤维采样一个长度 $l$，以构造出完整的纤维线段 $[\textbf a, \textbf b]$，其中 $\textbf a$ 和 $\textbf b$ 分别是纤维的起点和终点坐标。

纤维的长度 $l$ 直接影响着纤维在空间中的占用范围以及与其他纤维的交互情况。为了模拟实际 GDL 中纤维长度的分布特征，我们通常设定一个合理的长度范围 $[l_{\min}, l_{\max}]$，并从中采样一个长度值 $l$。具体地，我们可以使用一个均匀分布来采样长度：
\[
l \sim \mathcal{U}(l_{\min}, l_{\max})
\]
需要注意的是，在不重叠模式下，较长的纤维更容易与已有纤维发生碰撞，因此在实际实现中，我们可以根据当前的填充情况动态调整 $l_{\max}$，以缓解 jamming 问题。

对于每根纤维，在已知长度 $l$，中心点坐标 $\textbf c$ 和方向向量 $\textbf u$ 时，我们可以计算出纤维的起点 $\textbf a$ 和终点 $\textbf b$：
\[
\begin{cases}
\textbf a = \textbf c - \dfrac{l}{2} \textbf u \\
\\
\textbf b = \textbf c + \dfrac{l}{2} \textbf u
\end{cases}
\]
上述构造保证线段长度为 $|b-a|=l$，且线段中点为 $c$。该线段 $[a,b]$ 为候选中心轴线，后续将根据计算域边界执行裁剪（clipped）以得到位于域内的有效线段，详见第 2.2.5 节。

\subsubsection{\textbf{边界裁剪与体素化写入}}


